% Part: methods
% Chapter: proofs
% Section: starting-proofs

\documentclass[../../../include/open-logic-section]{subfiles}

\begin{document}

\olfileid{mth}{prf}{str}

\olsection{প্রমাণ শুরু করা}

কিন্তু শুরু করবে কোথা থেকে?

তোমাকে যে বক্তব্য প্রমাণ করতে দেওয়া হয়েছে, প্রমাণের
শেষে সেটিতেই পৌঁছতে হবে। অবশ্য শুরুতে \emph{ঘোষণা}
করতে পারো যে সেটি প্রমাণ করবে; কিন্তু একে অনুমান হিসেবে
ব্যবহার করা চলবে না। নতুন একটি কাগজের নিচে প্রমাণের
লক্ষ্যটি লিখে রাখো—তাহলে লক্ষ্য চোখের আড়াল হবে না।

এরপর হয়তো ব্যবহারের মতো কিছু অনুমান থাকবে (পরের বিভাগে
প্রমাণের \emph{ধরন} আলোচনা করলে বিষয়টি আরও স্পষ্ট হবে)।
সেগুলি কাগজের উপরে লেখো এবং যে এগুলি অনুমান, তা স্পষ্ট
করে দাও: যেমন, $p$ অনুমান করলে ``$p$ ধরি'' বা
``মনে করি, $p$'' লেখো। শেষে প্রশ্নে ব্যবহৃত কিছু সংজ্ঞাও
জানতে হবে। নির্দিষ্ট কোনো সংজ্ঞা ব্যবহার করতে বলা হতে
পারে, অথবা অনুমান বা যে সিদ্ধান্তে পৌঁছতে চাইছ, তাতেই
বিভিন্ন সংজ্ঞা থাকতে পারে। \emph{সেগুলি লিখে রাখো এবং
তাদের অর্থ বুঝেছ কি না নিশ্চিত হও।}

প্রমাণ সাজানোর রীতিও প্রশ্নের আকারের উপর নির্ভর করে।
পরের বিভাগে বাক্যের ধরন অনুযায়ী প্রমাণ সাজানোর বিস্তারিত
আলোচনা আছে।

\end{document}
